\practiceweek{7}

\begin{exercisebox}
\begin{pproblem}[$N = 8$, $a = 1$, and overshooting]\label{pp:grover8}
    Find the best $k$, and the success probability for $k = 1, 2, 3$.
\end{pproblem}
\tcblower
\begin{wsolution}
    \uses{\scr{cor:hk-1008}; \scr{thm:iterations-1008}}
    \leavevmode
    \begin{enumerate}
        \item $\sin\theta = \sfrac{1}{\sqrt{8}} \approx 0.3536$, so
              $\theta \approx 0.3614$ rad $\approx 20.7^\circ$.
        \item $k = \floor{\sfrac{\pi}{(4 \cdot 0.3614)}} = \floor{2.17} = 2$.
        \item Plug in:
    \end{enumerate}
    \begin{center}
        \begin{tabular}{ccc}
            \toprule
            $k$ & Angle $(2k+1)\theta$ & Success probability                \\ \midrule
            $0$ & $20.7^\circ$         & $\sfrac{1}{8} = 0.125$             \\
            $1$ & $62.1^\circ$         & $\sfrac{25}{32} \approx 0.781$     \\
            $2$ & $103.5^\circ$        & $\sfrac{121}{128} \approx 0.945$   \\
            $3$ & $144.9^\circ$        & $\sfrac{169}{512} \approx 0.330$   \\ \bottomrule
        \end{tabular}
    \end{center}
    The best is $k = 2$ with probability about $0.945$. A third iteration rotates past $90^\circ$
    and the probability drops to $0.33$: more iterations are not better.
\end{wsolution}
\end{exercisebox}

\begin{exercisebox}
\begin{pproblem}[Inversion about the mean for $N = 8$, $a = 1$]\label{wp:mean8}
    Track the eight amplitudes through one Grover iteration and compare with
    $\sin^2(3\theta)$.
\end{pproblem}
\tcblower
\begin{wsolution}
    \uses{\scr{lem:diffusion-1006}; \scr{cor:hk-1008}}
    \leavevmode
    \begin{enumerate}
        \item Start: every amplitude is $\sfrac{1}{\sqrt{8}}$.
        \item $Z_f$ negates the solution: seven amplitudes $\sfrac{1}{\sqrt{8}}$ and one
              $-\sfrac{1}{\sqrt{8}}$.
        \item The mean is $\frac{1}{8} \cdot \frac{7 - 1}{\sqrt{8}} = \frac{3}{4\sqrt{8}}$.
        \item Each amplitude $\alpha$ becomes $2 \cdot \text{mean} - \alpha$. Non-solutions:
              $\frac{6}{4\sqrt{8}} - \frac{1}{\sqrt{8}} = \frac{1}{2\sqrt{8}}$. Solution:
              $\frac{6}{4\sqrt{8}} + \frac{1}{\sqrt{8}} = \frac{5}{2\sqrt{8}}$.
        \item Success probability $\frac{25}{4 \cdot 8} = \frac{25}{32} \approx 0.781$. Check of the
              norm: $7 \cdot \frac{1}{32} + \frac{25}{32} = 1$.
    \end{enumerate}
    With $\sin\theta = \sfrac{1}{\sqrt{8}}$: $\sin 3\theta = 3\sin\theta - 4\sin^3\theta =
    \frac{3}{\sqrt{8}} - \frac{4}{8\sqrt{8}} = \frac{5}{2\sqrt{8}}$, the same amplitude.
\end{wsolution}
\end{exercisebox}

\begin{exercisebox}
\begin{pproblem}[$p = 1/100$]\label{pp:amp100}
    A quantum algorithm $\mathcal{A}$ outputs a good answer with probability $\sfrac{1}{100}$, and we
    can check an answer with $f$. How many runs of $\mathcal{A}$ are needed classically, and how
    many with amplitude amplification?
\end{pproblem}
\tcblower
\begin{wsolution}
    \uses{\scr{thm:ampamp-1008}}
    \leavevmode
    \begin{enumerate}
        \item Classically: run, check, repeat. The number of runs is geometric with mean
              $\sfrac{1}{p} = 100$.
        \item $\sin\theta = \sqrt{\sfrac{1}{100}} = 0.1$, so $\theta \approx 0.1002$.
        \item $k = \floor{\sfrac{\pi}{(4 \cdot 0.1002)}} = \floor{7.84} = 7$.
        \item The success probability is $\sin^2(15 \cdot 0.1002) = \sin^2(1.503) \approx 0.995$.
        \item Cost: $2 \cdot 7 + 1 = 15$ applications of $\mathcal{A}$ or its inverse, and $7$
              queries to $f$.
    \end{enumerate}
    So $15$ applications instead of $100$. In general the cost is at most
    $\frac{\pi}{2\sqrt{p}} + 1 = \bigO{\sfrac{1}{\sqrt{p}}}$ against $\sfrac{1}{p}$.
\end{wsolution}
\end{exercisebox}
